\section{Motivation}
Mobile endpoints continuously communicate across Wi-Fi, cellular, enterprise, and home networks,
yet defenders rarely have visibility into those flows without either rooting the device or
inspecting sensitive payload content. Android's \texttt{VpnService} provides a practical middle
ground by enabling user-space packet interception on non-rooted devices
\citep{android_vpnservice}. This makes metadata-only monitoring possible, but it also creates two
competing requirements: detection quality must be operationally useful, and data collection must
stay meaningfully privacy-constrained.

This tension is especially acute for encrypted traffic. Payload encryption removes content but not
all side channels. Flow timing, volume, burst shape, destination characteristics, and session
structure can still reveal application identity or user behavior. I therefore treat privacy as two
separate questions: what can be inferred from the representation MANTA releases, and what can be
inferred by a passive observer that sees the encrypted traffic on the wire.

\section{Problem Statement}
I study whether a practical, metadata-only IDS can be built and evaluated as an end-to-end artifact
while keeping the privacy boundary explicit. The project asks whether MANTA can:
\begin{enumerate}[leftmargin=*]
  \item capture and aggregate application-attributed flow metadata without payload inspection,
  \item support anomaly scoring across on-device, remote, privacy-student, and federated variants,
  \item quantify the privacy properties of released traffic representations, and
  \item distinguish those release-privacy properties from observer leakage visible in encrypted
  traffic sequences.
\end{enumerate}
The evaluation follows the same artifact through detection, release-view reduction, leakage
attacks, observer fingerprinting, and reproducibility packaging. That gives one place to compare
utility, privacy, and operating constraints instead of treating them as separate experiments.

\section{Research Questions}
\begin{itemize}[leftmargin=*]
  \item RQ1: To what extent can metadata-only flow windows support anomaly detection across a
  heterogeneous public encrypted-traffic corpus under grouped, source-aware evaluation?
  \item RQ2: How do on-device, remote, privacy-student, and federated variants compare in terms of
  anomaly-detection quality and practical deployment readiness?
  \item RQ3: What privacy-preserving properties do the reduced privacy views actually provide for
  \emph{released} traffic representations?
  \item RQ4: How different are those release-privacy properties from privacy against a passive
  observer of encrypted traffic sequences?
\end{itemize}

\section{Contributions}
\begin{itemize}[leftmargin=*]
  \item An end-to-end implementation that connects Android metadata capture, backend triage, and a
  reproducible machine-learning pipeline.
  \item A public-only training and evaluation corpus assembled from Android malware, enterprise,
  IoT, campus, and mobile-app traffic datasets with source/session-aware metadata carried through
  the entire pipeline.
  \item An evaluation protocol that compares on-device, remote, privacy-student, and
  federated variants under grouped source-aware holdouts instead of optimistic random splits.
  \item A split privacy methodology that distinguishes release privacy from observer leakage and
  measures both explicitly through metadata leakage attacks and encrypted-flow sequence audits.
  \item An empirical result showing that metadata coarsening can reduce exact app-ID leakage while
  still leaving substantial app-family, destination-behavior, and observer leakage.
\end{itemize}

\section{Thesis Structure}
Chapter 2 summarizes Android networking constraints, encrypted-traffic analysis, privacy
fingerprinting, and privacy-aware learning literature. Chapter 3 defines the system architecture,
trust boundaries, and the distinction between release privacy and observer privacy. Chapter 4
details the implementation of the Android endpoint, backend adapter, privacy views, privacy
student, and federated path. Chapter 5 presents the public-corpus evaluation methodology and
results. Chapter 6 discusses privacy-preserving properties, ethics, and limitations. Chapter 7
closes with future work.
